\section{Second Pass: Relations, Classification, and Quotients}
\chapterline

In Chapter~32, we used functions and fibers to decompose a domain into level sets. But that required an external target space and an existing map. Often we want to reshape a space from within: gluing boundaries together to form a torus, ignoring fine-grained distinctions, or grouping integers by remainder.

The formal tool for this is the \emph{equivalence relation}: a rule that formalizes which differences to track and which to collapse into a single object.

\begin{historybox}
In 1801, a 24-year-old Carl Friedrich Gauss published \emph{Disquisitiones Arithmeticae}, tackling a problem that had slowed number theory for generations: mathematicians treated division remainders ($a = qm + r$) as individual numerical values rather than mathematical objects with their own arithmetic laws. Gauss introduced modular congruence: $a \equiv b \pmod m$, meaning $a$ and $b$ leave the same remainder when divided by $m$, just like hours wrapping around a clock dial. Rather than computing with single integers, Gauss recognized that congruence bundles the infinite integers $\Z$ into a finite system of cyclical classes. Contemporaries questioned the notation, asking why one could not simply keep writing the equation $a - b = km$. Gauss proved that congruences behave like equality itself, allowing one to add and multiply whole classes of numbers: $[a] + [b] = [a+b]$.

A century later in 1920s G\"ottingen, Emmy Noether elevated equivalence classes into the central architecture of abstract algebra. Earlier nineteenth-century algebra relied heavily on explicit, coordinate-dependent polynomial computations. Leopold Kronecker resisted non-constructive sets and abstract classes (stating, ``God made the integers, all else is the work of man''), while Paul Gordan initially questioned abstract methods in invariant theory (calling Hilbert's non-constructive proof ``theology, not mathematics''). Noether developed a structural approach. She recognized that whenever a homomorphism $\phi$ connects algebraic systems, the kernel $\ker(\phi)$ (the collection of all inputs that map to the identity) induces an equivalence relation. Its quotient structure $G/\ker(\phi)$---the space formed by regarding each fiber as a single element---is isomorphic to the image. Her First Isomorphism Theorem, $G/\ker\phi \cong \operatorname{im}\phi$, eliminated lengthy coordinate computations in favor of clean structural classification.

Yet the general classification problem remains an active frontier today. In modern descriptive set theory (the study of definable infinite sets), Alexander Kechris and Greg Hjorth in the 1990s investigated when equivalence relations can be classified by concrete numerical invariants. They discovered the phenomenon of \emph{turbulence}: many natural equivalence relations (such as classifying orbits in chaotic dynamical systems, or classifying infinite-dimensional operator algebras) are provably so complex that no collection of numbers or discrete tags can ever classify their classes. The boundary between classifiable equivalence relations and turbulent ones remains an active area in mathematical logic.
\end{historybox}

\subsection{The Equivalence Machine: Reflexive, Symmetric, Transitive}

Before we can build equivalence classes, we need the general concept of a relation. A \emph{binary relation} $\sim$ on a set $A$ is simply a rule that inspects ordered pairs: given two elements $a, b \in A$, the relation declares either ``yes, $a$ is related to $b$'' (written $a \sim b$) or ``no, they are not'' ($a \not\sim b$). Formally, $\sim$ is a subset of the Cartesian product $A \times A$, consisting of all pairs $(a, b)$ that satisfy the condition.

A general binary relation need not impose any organizational structure on a set. To act as a genuine proxy for equality, a relation must satisfy three structural laws:
\begin{enumerate}
  \item \textbf{Reflexivity}: For all $a \in A$, $a \sim a$. (Every element is related to itself.)
  \item \textbf{Symmetry}: For all $a, b \in A$, if $a \sim b$, then $b \sim a$. (The relation is symmetric between pairs.)
  \item \textbf{Transitivity}: For all $a, b, c \in A$, if $a \sim b$ and $b \sim c$, then $a \sim c$. (Connections chain across intermediate elements.)
\end{enumerate}

To see why all three are indispensable, consider what happens when just one rule fails:
\begin{itemize}
  \item \textbf{Strict inequality ($<$) on $\R$}: Transitivity holds ($1 < 2$ and $2 < 5 \implies 1 < 5$), but reflexivity fails ($1 \not< 1$) and symmetry fails ($1 < 2$ but $2 \not< 1$).
  \item \textbf{Proximity (``within $1$ kilometer'')}: This relation is reflexive (you are within $0$ km of yourself) and symmetric (if Alice is near Bob, Bob is near Alice). But transitivity fails: Alice is $800$ meters from Bob, and Bob is $800$ meters from Charlie, yet Alice can easily be $1600$ meters from Charlie. Because transitivity fails, geographic proximity does not divide a population into disjoint subsets; instead, proximity neighborhoods overlap.
\end{itemize}

\begin{recoverybox}
\textbf{Mental check:} Let $A = \{1, 2, 3\}$. Consider the relation
\[
R = \{(1, 1), (2, 2), (3, 3), (1, 2), (2, 1)\} \subseteq A \times A.
\]
\begin{enumerate}[label=(\alph*)]
  \item Verify whether $R$ is reflexive, symmetric, and transitive.
  \item Find the equivalence class $[a] = \{x \in A : x \sim a\}$ for each element $a \in A$.
\end{enumerate}
\emph{Answer:} (a) Reflexivity holds because $(1, 1), (2, 2)$, and $(3, 3)$ all belong to $R$. Symmetry holds because the only off-diagonal pairs are $(1, 2)$ and $(2, 1)$, which mirror each other. For transitivity, chaining $(1, 2)$ and $(2, 1)$ requires $(1, 1) \in R$ (present), and chaining $(2, 1)$ and $(1, 2)$ requires $(2, 2) \in R$ (present). Thus $R$ is an equivalence relation. (b) For $1$, all elements related to $1$ are $[1] = \{1, 2\}$. Similarly, $[2] = \{1, 2\} = [1]$. For $3$, no off-diagonal pairs touch $3$, so $[3] = \{3\}$. Notice that $\{[1], [3]\} = \{\{1, 2\}, \{3\}\}$ partitions $\{1, 2, 3\}$ into two disjoint blocks.
\end{recoverybox}

These three conditions ensure that the set $A$ breaks apart cleanly into disjoint subsets called \emph{equivalence classes}:
\[
[a] = \{x \in A : x \sim a\}.
\]
Any element inside $[a]$ is called a \emph{representative} of the class.

\begin{figure}[htbp]
\centering
\begin{tikzpicture}[scale=0.95, >=Stealth]
  \tikzset{
    relcircle/.style={draw=sourceblue!80!black, fill=sourceblue!10, thick},
    relcircleb/.style={draw=targetgreen!80!black, fill=targetgreen!10, thick},
    pointdot/.style={circle, fill=sourceblue!90!black, inner sep=2pt},
    pointdotg/.style={circle, fill=targetgreen!90!black, inner sep=2pt},
    pointdotr/.style={circle, fill=trapred!90!black, inner sep=2.2pt},
    fusedbox/.style={draw=quotientorange!80!black, fill=quotientorange!10, very thick, rounded corners=8pt}
  }

  % LEFT PANEL: Proximity / General relation (Partial overlap allowed)
  \begin{scope}[shift={(0,0)}]
    \node[font=\small\bfseries, text=sourceblue!80!black] at (2.2, 2.4) {General Relation: Partial Overlap};
    
    % Two overlapping disks
    \draw[relcircle] (1.3, 0.4) circle (1.3);
    \draw[relcircleb] (3.1, 0.4) circle (1.3);
    
    % Elements
    \node[pointdot, label={above left:\small $a$}] (a) at (0.6, 0.4) {};
    \node[pointdotg, label={above right:\small $b$}] (b) at (3.8, 0.4) {};
    \node[pointdotr, label={above:\small $c$}] (c) at (2.2, 0.4) {};
    
    % Links
    \draw[thick, draw=sourceblue!80!black] (a) -- node[below, font=\scriptsize] {$a \sim c$} (c);
    \draw[thick, draw=targetgreen!80!black] (c) -- node[below, font=\scriptsize] {$c \sim b$} (b);
    \draw[thick, dashed, draw=trapred!80!black] (a) to[bend left=30] node[above, font=\scriptsize, text=trapred!90!black, fill=white, inner sep=1.5pt] {$a \nsim b$ (no transitivity)} (b);
    
    \node[below, font=\footnotesize, align=center, text=gray!80!black] at (2.2, -1.4) 
      {Sharing $c$ does not connect $a$ and $b$.\\The sets overlap partially without merging.};
  \end{scope}

  % DIVIDER
  \draw[thick, draw=gray!30, dashed] (5.2, -1.8) -- (5.2, 2.7);

  % RIGHT PANEL: Equivalence Relation (Transitivity forces fusion)
  \begin{scope}[shift={(5.9,0)}]
    \node[font=\small\bfseries, text=quotientorange!90!black] at (2.4, 2.4) {Equivalence Relation: Total Fusion};
    
    % Single unified boundary enclosing everything
    \draw[fusedbox] (0.0, -1.0) rectangle (4.8, 1.8);
    \node[font=\footnotesize\bfseries, text=quotientorange!90!black, anchor=north east] at (4.7, 1.7) {$[a] = [b]$};

    % Elements inside the fused class
    \node[pointdot, label={above left:\small $a$}] (ea) at (1.0, 0.4) {};
    \node[pointdotg, label={above right:\small $b$}] (eb) at (3.8, 0.4) {};
    \node[pointdotr, label={above:\small $c$}] (ec) at (2.4, 0.4) {};
    \node[pointdot, label={below left:\small $y$}] (ey) at (1.0, -0.4) {};

    % Transitivity chain
    \draw[thick, draw=sourceblue!80!black] (ea) -- node[above, font=\scriptsize] {$a \sim c$} (ec);
    \draw[thick, draw=targetgreen!80!black] (ec) -- node[above, font=\scriptsize] {$c \sim b$} (eb);
    \draw[thick, draw=quotientorange!90!black] (ey) -- node[left, font=\scriptsize] {$y \sim a$} (ea);
    \draw[thick, draw=quotientorange!90!black, -{Latex}] (ey) to[bend right=25] node[below, font=\scriptsize, text=quotientorange!90!black] {transitivity forces $y \sim b$} (eb);

    \node[below, font=\footnotesize, align=center, text=gray!80!black] at (2.4, -1.4) 
      {A single shared element $c$ chains $a \sim b$.\\Any $y \in [a]$ is pulled into $[b]$: $[a] \subseteq [b]$.};
  \end{scope}
\end{tikzpicture}
\caption{The Partition Principle. \textbf{Left}: In a general relation (such as metric distance $\le 1$), two relational neighborhoods can partially overlap: point $a$ is close to $c$, and $c$ is close to $b$, but $a$ is far from $b$. The overlap is a porous sliver. \textbf{Right}: In an equivalence relation, transitivity ensures that if two candidate classes $[a]$ and $[b]$ share even one element $c$, then $a \sim c$ and $c \sim b$ force $a \sim b$. Any element $y \in [a]$ is also in $[b]$, so $[a] \subseteq [b]$ and $[b] \subseteq [a]$. Two equivalence classes cannot partially overlap: they are either disjoint or completely identical.}
\label{33:fig-partition-principle}
\end{figure}

\begin{theorem}[The Fundamental Partition Theorem]
Let $\sim$ be an equivalence relation on a non-empty set $A$. Then the collection of equivalence classes $\{[a] : a \in A\}$ forms a partition of $A$:
\begin{enumerate}
  \item \textbf{Exhaustion}: Every element $x \in A$ belongs to at least one equivalence class ($x \in [x]$).
  \item \textbf{Disjointness}: Any two equivalence classes $[a]$ and $[b]$ are either completely identical or completely disjoint:
  \[
  [a] = [b] \quad \text{or} \quad [a] \cap [b] = \emptyset.
  \]
\end{enumerate}
Conversely, every partition of $A$ into non-empty disjoint blocks defines a unique equivalence relation whose classes are precisely the blocks of the partition.
\end{theorem}

\begin{proof}
Item (1) follows immediately from reflexivity: for any $x \in A$, $x \sim x$, so $x \in [x]$.

For item (2), suppose $[a] \cap [b] \neq \emptyset$. We must prove that $[a] = [b]$ identically.

Since the intersection is non-empty, choose an element $c \in [a] \cap [b]$. By definition of equivalence classes:
\[
c \sim a \quad \text{and} \quad c \sim b.
\]
By symmetry, $c \sim a$ implies $a \sim c$. Since $a \sim c$ and $c \sim b$, transitivity yields $a \sim b$.

Now let $y \in [a]$ be arbitrary. Then $y \sim a$. Because $y \sim a$ and $a \sim b$, transitivity guarantees that $y \sim b$, which means $y \in [b]$. Since $y$ was arbitrary, this establishes $[a] \subseteq [b]$.

By an identical argument swapping the roles of $a$ and $b$ (noting $b \sim a$ by symmetry), we obtain $[b] \subseteq [a]$. By mutual subset inclusion, $[a] = [b]$.

Finally, for the converse: given a partition $\mathcal{P} = \{B_i\}_{i \in I}$ of $A$, declare $x \sim y$ if and only if $x$ and $y$ lie in the same block $B_i$. Since each $x$ lies in a unique block, $x$ shares its block with itself (reflexive), sharing a block is mutual (symmetric), and if $x, y \in B_i$ and $y, z \in B_j$, uniqueness of the block containing $y$ forces $B_i = B_j$, so $x, z \in B_i$ (transitive). The equivalence classes are exactly the sets $B_i$.
\end{proof}

Figure~\ref{33:fig-partition-tiling} illustrates this correspondence: an equivalence relation partitions a set into non-overlapping subsets that cover the entire domain.

\begin{figure}[htbp]
\centering
\begin{tikzpicture}[scale=0.95, >=Latex]
  % Left Panel: Space partitioned into blocks
  \begin{scope}[xshift=0cm]
    \node[anchor=south, font=\small\bfseries, text=sourceblue!90!black] at (2.6, 2.3) {1. Tiling of Space $A$};
    \draw[rounded corners=6pt, fill=gray!3, draw=gray!30, thick] (-0.2, -2.1) rectangle (5.4, 2.3);
    
    % Block 1: [a]
    \draw[rounded corners=4pt, fill=sourceblue!15, draw=sourceblue!70!black, thick] (0.1, 0.3) rectangle (2.4, 1.9);
    \node[font=\footnotesize\bfseries, text=sourceblue!90!black] at (0.6, 1.6) {$[a]$};
    \fill[sourceblue!90!black] (0.7, 0.8) circle (2pt) node[below=2pt, font=\tiny] {$a$};
    \fill[sourceblue!90!black] (1.8, 1.2) circle (2pt) node[below=2pt, font=\tiny] {$x$};
    \draw[sourceblue!60!black, thin] (0.7, 0.8) -- (1.8, 1.2);
    
    % Block 2: [b]
    \draw[rounded corners=4pt, fill=targetgreen!15, draw=targetgreen!70!black, thick] (2.8, 0.3) rectangle (5.1, 1.9);
    \node[font=\footnotesize\bfseries, text=targetgreen!90!black] at (3.3, 1.6) {$[b]$};
    \fill[targetgreen!90!black] (3.4, 0.9) circle (2pt) node[below=2pt, font=\tiny] {$b$};
    \fill[targetgreen!90!black] (4.5, 1.1) circle (2pt) node[below=2pt, font=\tiny] {$w$};
    \draw[targetgreen!60!black, thin] (3.4, 0.9) -- (4.5, 1.1);
    
    % Block 3: [c]
    \draw[rounded corners=4pt, fill=quotientorange!15, draw=quotientorange!70!black, thick] (0.1, -1.3) rectangle (2.4, -0.1);
    \node[font=\footnotesize\bfseries, text=quotientorange!90!black] at (0.6, -0.35) {$[c]$};
    \fill[quotientorange!90!black] (0.8, -0.8) circle (2pt) node[below=2pt, font=\tiny] {$c$};
    \fill[quotientorange!90!black] (1.8, -0.7) circle (2pt) node[below=2pt, font=\tiny] {$z$};
    \draw[quotientorange!60!black, thin] (0.8, -0.8) -- (1.8, -0.7);
    
    % Block 4: [d]
    \draw[rounded corners=4pt, fill=structurepurple!15, draw=structurepurple!70!black, thick] (2.8, -1.3) rectangle (5.1, -0.1);
    \node[font=\footnotesize\bfseries, text=structurepurple!90!black] at (3.3, -0.35) {$[d]$};
    \fill[structurepurple!90!black] (3.5, -0.8) circle (2pt) node[below=2pt, font=\tiny] {$d$};
    \fill[structurepurple!90!black] (4.4, -0.7) circle (2pt) node[below=2pt, font=\tiny] {$u$};
    \draw[structurepurple!60!black, thin] (3.5, -0.8) -- (4.4, -0.7);
    
    \node[font=\scriptsize, align=center, text=gray!80!black] at (2.6, -1.75) {Tiling: $\bigcup [x] = A$ and $[x] \cap [y] = \emptyset$};
  \end{scope}

  % Right Panel: Dichotomy Theorem
  \begin{scope}[xshift=6.4cm]
    \node[anchor=south, font=\small\bfseries, text=targetgreen!90!black] at (2.6, 2.3) {2. Disjoint or Identical};
    \draw[rounded corners=6pt, fill=gray!3, draw=gray!30, thick] (-0.2, -2.1) rectangle (5.4, 2.3);
    
    % Case 1: a not ~ b (Disjoint)
    \node[font=\scriptsize\bfseries, text=sourceblue!90!black] at (2.6, 1.7) {Case 1: $a \nsim b \implies [a] \cap [b] = \emptyset$ (Disjoint)};
    \draw[rounded corners=3pt, fill=sourceblue!12, draw=sourceblue!70!black, thick] (0.3, 0.4) rectangle (2.2, 1.4);
    \node[font=\tiny\bfseries, text=sourceblue!90!black] at (0.7, 1.1) {$[a]$};
    \fill[sourceblue!90!black] (1.3, 0.8) circle (2pt) node[right=2pt, font=\tiny] {$a$};
    
    \draw[rounded corners=3pt, fill=targetgreen!12, draw=targetgreen!70!black, thick] (3.0, 0.4) rectangle (4.9, 1.4);
    \node[font=\tiny\bfseries, text=targetgreen!90!black] at (3.4, 1.1) {$[b]$};
    \fill[targetgreen!90!black] (4.0, 0.8) circle (2pt) node[right=2pt, font=\tiny] {$b$};
    
    \draw[thick, gray!40, dashed] (0.2, 0.0) -- (5.0, 0.0);
    
    % Case 2: a ~ b (Identical)
    \node[font=\scriptsize\bfseries, text=quotientorange!90!black] at (2.6, -0.3) {Case 2: $a \sim b \implies [a] = [b]$ (Total Fusion)};
    \draw[rounded corners=4pt, fill=quotientorange!15, draw=quotientorange!80!black, thick] (0.8, -1.3) rectangle (4.4, -0.6);
    \node[font=\tiny\bfseries, text=quotientorange!90!black] at (1.4, -0.85) {$[a] = [b]$};
    \fill[quotientorange!90!black] (2.7, -0.95) circle (2pt) node[below=1pt, font=\tiny] {$a$};
    \fill[quotientorange!90!black] (3.7, -0.95) circle (2pt) node[below=1pt, font=\tiny] {$b$};
    \draw[thick, quotientorange!80!black] (2.7, -0.95) -- node[above, font=\tiny] {$\sim$} (3.7, -0.95);
    
    \node[font=\scriptsize, align=center, text=gray!80!black] at (2.6, -1.75) {Transitivity eliminates partial overlap};
  \end{scope}
\end{tikzpicture}
\caption{The Fundamental Partition Theorem. \textbf{Left}: An equivalence relation partitions the domain $A$ into non-overlapping subsets that cover every element. \textbf{Right}: For any two elements $a$ and $b$, their equivalence classes are either completely disjoint (when $a \nsim b$) or completely identical (when $a \sim b$). No partial overlap can exist.}
\label{33:fig-partition-tiling}
\end{figure}

\subsection{The Well-Definedness Trap and Wild Quotients}

Passing from individual elements to equivalence classes introduces two classic hazards:

\begin{itemize}
  \item \textbf{The Well-Definedness Criterion}: Formulas on a quotient set $X/\!\sim$ are almost always written by picking a representative $a$ of the class $[a]$ and applying some rule $f(a)$. Because a single class $[a]$ has many different names, you have to check that picking another representative $a' \sim a$ produces the exact same result:
  \[
  a \sim a' \implies f(a) = f(a').
  \]
  If that check fails, the proposed rule $F([a]) = f(a)$ depends on the name you picked rather than the class itself, so it is not \emph{well-defined} and does not define a function at all. For example, in the construction of the rational numbers $\Q$ as pairs of integers $[(p, q)]$ with $q \neq 0$ under $(p, q) \sim (r, s) \iff ps = qr$, attempting to define the ``numerator'' by $N([(p, q)]) = p$ is invalid: $[(1, 2)] = [(2, 4)]$, but $N([(1, 2)]) = 1 \neq 2 = N([(2, 4)])$. In contrast, addition $[(p, q)] + [(r, s)] = [(ps + qr, qs)]$ is well-defined because cross-multiplication confirms that replacing $(p, q)$ or $(r, s)$ with equivalent pairs leaves the resulting class unchanged.

  Similarly, in the integers modulo $4$ ($\Z/4\Z$), defining a function $f: \Z/4\Z \to \Z/2\Z$ by $f([x]_4) = [x]_2$ is valid, because $x \equiv x' \pmod 4 \implies 4 \mid (x - x') \implies 2 \mid (x - x')$, so $[x]_2 = [x']_2$. But attempting to define $g: \Z/4\Z \to \Z/4\Z$ by $g([x]_4) = [2^x]_4$ is not well-defined: $[1]_4 = [5]_4$, yet $2^1 = 2 \not\equiv 2^5 = 32 \equiv 0 \pmod 4$.
  \item \textbf{Quotient Topologies and Separation Axioms}: Gluing points together in a geometric space can destroy nice separation properties even when the starting space is completely ordinary.
  
  For example, consider the real line $\R$ and identify any two points that differ by a rational number: $x \sim y \iff x - y \in \Q$. Each equivalence class is a coset $x + \Q$, which is a dense subset of the real line. In the quotient space $\R/\Q$, every non-empty open set contains points from every equivalence class. As a result, the quotient topology on $\R/\Q$ is the indiscrete topology: the only open sets are the empty set and the entire space. No two distinct points can be separated by disjoint open neighborhoods, so the Hausdorff property fails completely.
\end{itemize}

\begin{recoverybox}
\textbf{Mental check:} Consider $\Z/6\Z$, the integers modulo $6$. We want to define a map $H: \Z/6\Z \to \Z/m\Z$ using the rule $H([x]_6) = [x]_m$. For which values of $m \in \{2, 3, 4\}$ is $H$ well-defined?
\begin{enumerate}[label=(\alph*)]
  \item Test $m = 3$.
  \item Test $m = 4$.
\end{enumerate}
\emph{Answer:} (a) For $m = 3$, suppose $[x]_6 = [x']_6$. Then $x - x' = 6k = 3(2k)$ for some integer $k$. Thus $x \equiv x' \pmod 3$, which means $[x]_3 = [x']_3$. The map is well-defined. (b) For $m = 4$, in $\Z/6\Z$ we have $[0]_6 = [6]_6$. But $H([0]_6) = [0]_4$, while $H([6]_6) = [6]_4 = [2]_4$. Since $[0]_4 \neq [2]_4$, the rule produces conflicting outputs for the same class. In general, $H([x]_n) = [x]_m$ is well-defined if and only if $m$ divides $n$.
\end{recoverybox}

\subsection{Quotients: Turning Classes into Points}

Once we have partitioned a set $A$ into equivalence classes, we can treat each equivalence class as an element in its own right in a new space. 

This new space is the \emph{quotient set}, denoted $A/\sim$:
\[
A/\sim \;=\; \{[a] : a \in A\}.
\]
Accompanying this set is the \emph{canonical quotient map} $\pi: A \to A/\sim$, defined by $\pi(x) = [x]$, which sends each original element to its home class. The word ``canonical'' means that the map requires no arbitrary choices; the equivalence relation itself dictates where each point must go.

Figure~\ref{33:fig-clock-modular} visualizes this transformation in the familiar setting of modular arithmetic. The integers $\mathbb{Z}$ form an infinite line partitioned into four modulo-$4$ congruence classes. The canonical projection $\pi: \mathbb{Z} \to \mathbb{Z}/4\mathbb{Z}$ wraps the infinite line around a four-position clock dial, collapsing the infinite fiber $\{\dots, -4, 0, 4, 8, \dots\}$ into the single point $[0]$, and similarly for each remainder class.

\begin{figure}[htbp]
\centering
\begin{tikzpicture}[scale=0.92, >=Latex]
  % Left panel: The Infinite Number Line Z
  \begin{scope}[xshift=0cm]
    \node[anchor=south, font=\small\bfseries, text=sourceblue!90!black] at (2.7, 2.5) {1. Infinite Line $\mathbb{Z}$ Partitioned};
    
    % Axis
    \draw[thick, gray!60, <->] (-0.8, 0) -- (6.2, 0);
    
    % Colors:
    % 0 mod 4: targetgreen!80!black
    % 1 mod 4: sourceblue!80!black
    % 2 mod 4: quotientorange!80!black
    % 3 mod 4: structurepurple!80!black
    
    \foreach \n/\x/\c/\lbl in {-1/-0.5/structurepurple!80!black/{[3]}, 0/0.5/targetgreen!80!black/{[0]}, 1/1.5/sourceblue!80!black/{[1]}, 2/2.5/quotientorange!80!black/{[2]}, 3/3.5/structurepurple!80!black/{[3]}, 4/4.5/targetgreen!80!black/{[0]}, 5/5.5/sourceblue!80!black/{[1]}} {
      \draw[thick, gray!50] (\x, -0.1) -- (\x, 0.1);
      \fill[\c] (\x, 0) circle (2.5pt);
      \node[above=2pt, font=\scriptsize\bfseries] at (\x, 0.1) {$\n$};
      \node[below=2pt, font=\tiny, text=\c] at (\x, -0.1) {$\lbl$};
    }
    
    % Brackets or periodic repeat arcs
    \draw[targetgreen!70!black, thick, dashed, ->] (0.5, 0.8) to[bend left=45] node[above, font=\tiny] {$+4$} (4.5, 0.8);
    \draw[sourceblue!70!black, thick, dashed, ->] (1.5, 0.6) to[bend left=45] node[above, font=\tiny] {$+4$} (5.5, 0.6);
    
    \node[align=center, font=\scriptsize, text=gray!80!black] at (2.7, -1.2) {Equivalence: $x \sim y \iff x \equiv y \pmod 4$\\Each fiber $k + 4\mathbb{Z}$ has infinite size};
  \end{scope}

  % Middle transition arrow
  \draw[->, very thick, structurepurple!80!black] (6.0, 0.6) -- node[above=2pt, font=\footnotesize\bfseries] {$\pi$} node[below=2pt, font=\tiny, align=center] {wrap line\\into dial} (7.2, 0.6);

  % Right panel: The Clock Dial Z/4Z
  \begin{scope}[xshift=10.8cm, yshift=0.6cm]
    \node[anchor=south, font=\small\bfseries, text=targetgreen!90!black] at (0, 1.9) {2. Quotient Ring $\mathbb{Z}/4\mathbb{Z}$};
    
    % Circle
    \draw[thick, gray!30] (0, 0) circle (1.4);
    
    % Clock nodes at 0, 1, 2, 3
    % 0 at top (0, 1.4)
    \fill[targetgreen!80!black] (0, 1.4) circle (3pt);
    \node[above=3pt, font=\scriptsize\bfseries, text=targetgreen!90!black] at (0, 1.4) {$[0] = 4\mathbb{Z}$};
    
    % 1 at right (1.4, 0)
    \fill[sourceblue!80!black] (1.4, 0) circle (3pt);
    \node[right=3pt, font=\scriptsize\bfseries, text=sourceblue!90!black] at (1.4, 0) {$[1] = 1+4\mathbb{Z}$};
    
    % 2 at bottom (0, -1.4)
    \fill[quotientorange!80!black] (0, -1.4) circle (3pt);
    \node[below=3pt, font=\scriptsize\bfseries, text=quotientorange!90!black] at (0, -1.4) {$[2] = 2+4\mathbb{Z}$};
    
    % 3 at left (-1.4, 0)
    \fill[structurepurple!80!black] (-1.4, 0) circle (3pt);
    \node[left=3pt, font=\scriptsize\bfseries, text=structurepurple!90!black] at (-1.4, 0) {$[3] = 3+4\mathbb{Z}$};
    
    % Circular flow arrows
    \draw[->, thick, gray!60] (0.3, 1.35) arc (75:15:1.4);
    \draw[->, thick, gray!60] (1.35, -0.3) arc (-15:-75:1.4);
    \draw[->, thick, gray!60] (-0.3, -1.35) arc (-105:-165:1.4);
    \draw[->, thick, gray!60] (-1.35, 0.3) arc (165:105:1.4);
    
    \node[align=center, font=\scriptsize, text=gray!80!black] at (0, -2.4) {Infinite fibers collapse into\\4 distinct quotient points};
  \end{scope}
\end{tikzpicture}
\caption{The quotient map as periodic wrapping. \textbf{Left}: The integer line $\mathbb{Z}$ partitioned into four modulo-$4$ arithmetic progressions (color-coded). \textbf{Right}: The canonical projection $\pi(n) = [n]_4$ wraps the infinite line into a $4$-element cyclic clock $\mathbb{Z}/4\mathbb{Z}$. Each single point on the clock represents an entire infinite fiber.}
\label{33:fig-clock-modular}
\end{figure}

\begin{figure}[htbp]
\centering
\begin{tikzpicture}[scale=1.0, >=Stealth]
  \tikzset{
    domainbox/.style={draw=sourceblue!80!black, fill=sourceblue!6, very thick, rounded corners=4pt},
    classbox/.style={draw=sourceblue!60!black, fill=sourceblue!15, thick, rounded corners=3pt, inner sep=3pt},
    pointdot/.style={circle, fill=sourceblue!80!black, inner sep=1.8pt},
    quotbox/.style={draw=quotientorange!80!black, fill=quotientorange!8, very thick, rounded corners=4pt},
    quotpoint/.style={circle, fill=quotientorange!90!black, inner sep=2.5pt},
    targbox/.style={draw=targetgreen!80!black, fill=targetgreen!8, very thick, rounded corners=4pt},
    targpoint/.style={circle, fill=targetgreen!90!black, inner sep=2.5pt},
    arrowstyle/.style={->, thick, shorten >=3pt, shorten <=3pt}
  }

  % PANEL 1: Original Space X partitioned into equivalence classes
  \node[font=\footnotesize\scshape, text=sourceblue!80!black] at (1.6, 4.15) {Domain Space $X$ (Fibers)};
  \draw[domainbox] (0, 0) rectangle (3.2, 3.8);

  % Equivalence class regions in X
  \draw[classbox] (0.2, 2.6) rectangle (3.0, 3.6);
  \node[pointdot, label={above:\tiny $x_1$}] (p1) at (0.8, 2.95) {};
  \node[pointdot, label={above:\tiny $x_2$}] (p2) at (1.9, 2.95) {};
  \node[font=\scriptsize\bfseries, text=sourceblue!90!black, anchor=east] at (2.9, 3.35) {$[x_1]$};

  \draw[classbox] (0.2, 1.4) rectangle (3.0, 2.4);
  \node[pointdot, label={above:\tiny $x_3$}] (p3) at (0.7, 1.75) {};
  \node[pointdot, label={above:\tiny $x_4$}] (p4) at (1.4, 1.75) {};
  \node[pointdot, label={above:\tiny $x_5$}] (p5) at (2.1, 1.75) {};
  \node[font=\scriptsize\bfseries, text=sourceblue!90!black, anchor=east] at (2.9, 2.15) {$[x_3]$};

  \draw[classbox] (0.2, 0.2) rectangle (3.0, 1.2);
  \node[pointdot, label={above:\tiny $x_6$}] (p6) at (0.9, 0.55) {};
  \node[pointdot, label={above:\tiny $x_7$}] (p7) at (2.0, 0.55) {};
  \node[font=\scriptsize\bfseries, text=sourceblue!90!black, anchor=east] at (2.9, 0.95) {$[x_6]$};

  % PANEL 2: Quotient Space X / ~
  \draw[quotbox] (0, -4.2) rectangle (3.2, -1.5);
  \node[font=\footnotesize\scshape, text=quotientorange!90!black] at (1.6, -1.8) {Quotient Space $X/\sim$};
  \node[quotpoint, label={right:\footnotesize $[x_1] = \{x_1, x_2\}$}] (qa) at (0.8, -2.4) {};
  \node[quotpoint, label={right:\footnotesize $[x_3] = \{x_3, x_4, x_5\}$}] (qb) at (0.8, -3.1) {};
  \node[quotpoint, label={right:\footnotesize $[x_6] = \{x_6, x_7\}$}] (qc) at (0.8, -3.8) {};

  % PANEL 3: Target Space Y & Image im(f)
  \node[font=\footnotesize\scshape, text=targetgreen!80!black] at (8.0, 4.15) {Target Space $Y$};
  \draw[targbox] (6.4, 0) rectangle (9.6, 3.8);
  \node[font=\scriptsize\scshape, text=targetgreen!80!black] at (8.0, 3.55) {$\operatorname{im}(f) \subseteq Y$};
  \draw[draw=targetgreen!60!black, fill=targetgreen!15, thick, rounded corners=3pt] (6.6, 0.4) rectangle (9.4, 3.3);
  \node[targpoint, label={right:\footnotesize $y_1 = f(x_1) = f(x_2)$}] (ya) at (7.0, 2.7) {};
  \node[targpoint, label={right:\footnotesize $y_2 = f(x_3) = \dots$}] (yb) at (7.0, 1.8) {};
  \node[targpoint, label={right:\footnotesize $y_3 = f(x_6) = f(x_7)$}] (yc) at (7.0, 0.9) {};

  % Connecting Arrows: First Isomorphism Theorem
  % f: X -> Y
  \draw[arrowstyle, draw=sourceblue!80!black] (3.3, 2.5) -- node[above, font=\small] {$f$} (6.3, 2.5);

  % \pi: X -> X/~
  \draw[arrowstyle, draw=structurepurple!80!black] (1.6, -0.05) -- node[left=3pt, font=\footnotesize] {$\pi$ (projection: glue)} (1.6, -1.45);

  % \bar{f}: X/~ -> im(f)
  \draw[arrowstyle, draw=quotientorange!90!black, dashed] (3.3, -2.8) to[bend right=20] node[below right=2pt, font=\small] {$\bar{f}$ (unique bijection)} (6.7, 0.5);

  % Commutativity note
  \node at (4.8, 1.0) [font=\footnotesize, text=gray!70!black, align=center] {Commutative Diagram:\\$f = \bar{f} \circ \pi$};
\end{tikzpicture}
\caption{The First Isomorphism Theorem architecture. Any function $f: X \to Y$ partitions its domain into fibers of elements that have the same output value. Passing to the quotient $X/\sim$ under the fiber relation $x \sim x' \iff f(x) = f(x')$ identifies each fiber with a single element via the canonical projection $\pi(x) = [x]$. The induced map $\bar{f}([x]) = f(x)$ is well-defined, injective, and surjective onto $\operatorname{im}(f)$, establishing a canonical bijection $X/\sim \;\cong\; \operatorname{im}(f)$ that factors $f$ as $\bar{f} \circ \pi$.}
\label{33:fig-first-isomorphism}
\end{figure}

\subsection{The First Isomorphism Theorem}

Notice the structural insight that Figure~\ref{33:fig-first-isomorphism} reveals. Any function $f: X \to Y$ induces an equivalence relation on its domain:
\[
x \sim x' \iff f(x) = f(x').
\]
The equivalence classes $[x]$ are precisely the \emph{fibers} of the function---the inverse images $f^{-1}(\{y\})$ of individual points in the target space. 

The canonical projection map $\pi: X \to X/\sim$ sends each element to its equivalence class. The induced map $\bar{f}$ then connects the quotient directly to the image:

\begin{theorem}[The First Isomorphism Theorem for Sets]
Let $f: X \to Y$ be a function. Define the fiber equivalence relation on $X$ by:
\[
x \sim x' \iff f(x) = f(x').
\]
Let $\pi: X \to X/\sim$ be the canonical quotient projection, $\pi(x) = [x]$. Then:
\begin{enumerate}
  \item There exists a unique function $\bar{f}: X/\sim \;\to \operatorname{im}(f)$ such that $f = \bar{f} \circ \pi$.
  \item The induced function $\bar{f}$ is a bijection, establishing an isomorphism of sets:
  \[
  X/\sim \;\cong\; \operatorname{im}(f).
  \]
\end{enumerate}
\end{theorem}

\begin{proof}
To satisfy $f = \bar{f} \circ \pi$, we must have $f(x) = \bar{f}(\pi(x)) = \bar{f}([x])$ for every $x \in X$. This completely determines our formula: we are forced to define $\bar{f}([x]) = f(x)$.

We now verify the three essential properties:
\begin{enumerate}
  \item \textbf{Well-definedness}: If $[x] = [x']$, then $x \sim x'$. By definition of the relation, this means $f(x) = f(x')$. Hence $\bar{f}([x]) = \bar{f}([x'])$, so the output does not depend on which representative we picked. This proves existence and uniqueness of the map $\bar{f}$.
  \item \textbf{Injectivity}: Suppose $\bar{f}([x]) = \bar{f}([x'])$. By definition of $\bar{f}$, this means $f(x) = f(x')$. By definition of $\sim$, this gives $x \sim x'$, which implies $[x] = [x']$. The map never conflates distinct classes.
  \item \textbf{Surjectivity}: Let $y \in \operatorname{im}(f)$ be arbitrary. By definition of the image, there exists some $x \in X$ such that $f(x) = y$. Then $\bar{f}([x]) = f(x) = y$. Every element of the image is hit.
\end{enumerate}
Therefore, $\bar{f}$ is a well-defined bijection between $X/\sim$ and $\operatorname{im}(f)$.
\end{proof}

In abstract algebra, where $f: G \to H$ is a group homomorphism, the fiber over the identity $e_H$ is the kernel $\ker(f)$. The equivalence classes are the cosets $g\ker(f)$, and $\bar{f}$ automatically preserves the group operation:
\[
\bar{f}([a][b]) = \bar{f}([ab]) = f(ab) = f(a)f(b) = \bar{f}([a])\bar{f}([b]).
\]
The set bijection $\bar{f}$ yields a group isomorphism: $G/\ker(f) \cong \operatorname{im}(f)$.

To see the First Isomorphism Theorem geometrically before turning to boundary gluing, consider the squared-distance function $f: \mathbb{R}^2 \to [0, \infty)$ defined by $f(x, y) = x^2 + y^2$. Figure~\ref{33:fig-concentric-circles} illustrates how the fiber equivalence relation decomposes the plane into concentric circles centered at the origin. The canonical quotient projection $\pi$ identifies each circle with a single element, while the induced bijection $\bar{f}$ pairs each circle with its unique radius squared, yielding a quotient space identified with the half-line $[0, \infty)$.

\begin{figure}[htbp]
\centering
\begin{tikzpicture}[scale=0.95, >=Latex]
  % Left panel: 2D Plane with Concentric Circular Fibers
  \begin{scope}[xshift=0cm]
    \node[anchor=south, font=\small\bfseries, text=sourceblue!90!black] at (0, 2.6) {1. Plane $\mathbb{R}^2$ with Fibers};
    
    % Coordinate axes
    \draw[gray!40, ->] (-2.2, 0) -- (2.2, 0) node[right, font=\tiny] {$x$};
    \draw[gray!40, ->] (0, -2.1) -- (0, 2.1) node[right, font=\tiny] {$y$};
    
    % Origin fiber (orange)
    \fill[quotientorange!90!black] (0, 0) circle (2.5pt);
    \node[below left=1pt, font=\tiny, text=quotientorange!90!black] at (0, 0) {$\mathbf{0}$};
    
    % Concentric circles
    % r = 0.7
    \draw[thick, sourceblue!80!black] (0, 0) circle (0.7);
    \fill[sourceblue!80!black] (0.49, 0.49) circle (2pt) node[above right=-1pt, font=\tiny] {$p$};
    \fill[sourceblue!80!black] (-0.49, 0.49) circle (2pt) node[above left=-1pt, font=\tiny] {$q$};
    \node[font=\tiny\bfseries, text=sourceblue!90!black] at (0.8, -0.45) {$C_{1}$};
    
    % r = 1.3
    \draw[thick, targetgreen!80!black] (0, 0) circle (1.3);
    \node[font=\tiny\bfseries, text=targetgreen!90!black] at (1.4, -0.65) {$C_{2}$};
    
    % r = 1.9
    \draw[thick, structurepurple!80!black] (0, 0) circle (1.9);
    \node[font=\tiny\bfseries, text=structurepurple!90!black] at (1.95, -0.85) {$C_{3}$};
    
    \node[font=\scriptsize, align=center, text=gray!80!black] at (0, -2.6) {Fibers: $f^{-1}(\{c\}) = \{p : \|p\|^2 = c\}$\\$p \sim q \iff f(p) = f(q)$};
  \end{scope}

  % Middle Arrows
  \begin{scope}[xshift=3.0cm]
    \draw[->, very thick, structurepurple!80!black] (0, 0.6) -- node[above=2pt, font=\footnotesize\bfseries] {$\pi$} node[below=2pt, font=\tiny, align=center] {collapse\\circles} (1.4, 0.6);
    \draw[->, very thick, quotientorange!90!black, dashed] (0, -0.6) -- node[above=2pt, font=\footnotesize\bfseries] {$\bar{f}$} node[below=2pt, font=\tiny, align=center] {bijection} (1.4, -0.6);
  \end{scope}

  % Right panel: Target Ray [0, \infty)
  \begin{scope}[xshift=5.8cm]
    \node[anchor=south, font=\small\bfseries, text=targetgreen!90!black] at (2.4, 2.6) {2. Target Ray $[0, \infty)$ vs $\mathbb{R}^2/\sim$};
    
    % Half-line axis
    \draw[thick, gray!50, ->] (-0.2, 0) -- (4.8, 0) node[right, font=\tiny] {$c$};
    
    % Points on ray
    \fill[quotientorange!90!black] (0, 0) circle (3pt) node[below=5pt, font=\scriptsize\bfseries] {$0$};
    \node[above=5pt, font=\tiny, text=quotientorange!90!black] at (0, 0) {Origin};
    
    \fill[sourceblue!80!black] (1.1, 0) circle (3pt) node[below=5pt, font=\scriptsize\bfseries] {$c_1$};
    \node[above=5pt, font=\tiny, text=sourceblue!90!black] at (1.1, 0) {Fiber $C_1$};
    
    \fill[targetgreen!80!black] (2.3, 0) circle (3pt) node[below=5pt, font=\scriptsize\bfseries] {$c_2$};
    \node[above=5pt, font=\tiny, text=targetgreen!90!black] at (2.3, 0) {Fiber $C_2$};
    
    \fill[structurepurple!80!black] (3.7, 0) circle (3pt) node[below=5pt, font=\scriptsize\bfseries] {$c_3$};
    \node[above=5pt, font=\tiny, text=structurepurple!90!black] at (3.7, 0) {Fiber $C_3$};
    
    \node[font=\scriptsize, align=center, text=gray!80!black] at (2.4, -2.6) {Each 1D circular fiber becomes 1 point;\\Quotient $\mathbb{R}^2/\sim$ is homeomorphic to $[0, \infty)$};
  \end{scope}
\end{tikzpicture}
\caption{Geometric fiber decomposition under the First Isomorphism Theorem. The function $f(x, y) = x^2 + y^2$ sends $\mathbb{R}^2$ to $[0, \infty)$. \textbf{Left}: The fibers $f^{-1}(\{c\})$ are concentric circles of radius $\sqrt{c}$ (and the origin for $c=0$). Two points $p, q$ are equivalent if and only if they lie on the same circle. \textbf{Right}: The canonical projection $\pi$ collapses each entire circle into a single point, yielding the quotient $\mathbb{R}^2/\sim$, which the induced map $\bar{f}$ bijectively identifies with the non-negative ray $[0, \infty)$.}
\label{33:fig-concentric-circles}
\end{figure}

\subsection{Geometric Gluing: Quotients in Topology}

Quotients are also a primary construction in geometry and topology. Identifying boundary points or opposite edges corresponds formally to defining an equivalence relation and passing to the quotient space.

\begin{figure}[htbp]
\centering
\begin{tikzpicture}[scale=0.9, >=Stealth]
  \tikzset{
    intervalstyle/.style={very thick, draw=sourceblue!80!black},
    edgestylea/.style={very thick, draw=sourceblue!80!black},
    edgestyleb/.style={very thick, draw=targetgreen!80!black},
    edgestylerev/.style={very thick, draw=trapred!80!black},
    pointdot/.style={circle, fill=sourceblue!90!black, inner sep=2pt}
  }

  % PANEL 1: Interval [0, 1] to Circle S^1
  \begin{scope}[shift={(0,0)}]
    \node[font=\small\bfseries, text=sourceblue!80!black] at (1.6, 2.6) {Circle $S^1 = [0, 1]/(0 \sim 1)$};
    
    % Interval
    \draw[intervalstyle] (0.2, 1.6) -- (3.0, 1.6);
    \node[pointdot, label={above:\small $0$}] at (0.2, 1.6) {};
    \node[pointdot, label={above:\small $1$}] at (3.0, 1.6) {};
    \node[font=\footnotesize, text=gray!70!black] at (1.6, 1.2) {Unit interval $[0, 1]$};
    
    % Bending arrow
    \draw[->, thick, draw=structurepurple!80!black] (1.6, 0.9) -- node[right, font=\tiny] {glue $0 \sim 1$} (1.6, 0.1);
    
    % Circle
    \draw[very thick, draw=targetgreen!80!black] (1.6, -0.9) circle (0.8);
    \node[circle, fill=targetgreen!90!black, inner sep=2pt, label={right:\footnotesize $[0] = [1]$}] at (2.4, -0.9) {};
  \end{scope}

  % PANEL 2: Torus T^2 from Square
  \begin{scope}[shift={(4.5,0)}]
    \node[font=\small\bfseries, text=targetgreen!80!black] at (1.6, 2.6) {Torus $T^2$: Flat Torus};
    
    % Square
    \fill[sourceblue!5] (0.2, -0.8) rectangle (3.0, 2.0);
    \draw[gray!30, thick] (0.2, -0.8) rectangle (3.0, 2.0);
    
    % Horizontal edges (single arrow, matching right)
    \draw[edgestylea, -{Latex[length=3mm]}] (0.2, -0.8) -- (1.8, -0.8);
    \draw[edgestylea] (1.8, -0.8) -- (3.0, -0.8);
    \draw[edgestylea, -{Latex[length=3mm]}] (0.2, 2.0) -- (1.8, 2.0);
    \draw[edgestylea] (1.8, 2.0) -- (3.0, 2.0);
    \node[below, font=\scriptsize, text=sourceblue!80!black] at (1.6, -0.8) {$(x, 0) \sim (x, 1)$};
    
    % Vertical edges (double arrow, matching up)
    \draw[edgestyleb, -{Latex[length=2.5mm]}] (0.2, -0.8) -- (0.2, 0.6);
    \draw[edgestyleb, -{Latex[length=2.5mm]}] (0.2, 0.4) -- (0.2, 1.2);
    \draw[edgestyleb] (0.2, 1.0) -- (0.2, 2.0);
    
    \draw[edgestyleb, -{Latex[length=2.5mm]}] (3.0, -0.8) -- (3.0, 0.6);
    \draw[edgestyleb, -{Latex[length=2.5mm]}] (3.0, 0.4) -- (3.0, 1.2);
    \draw[edgestyleb] (3.0, 1.0) -- (3.0, 2.0);
    \node[right, font=\scriptsize, text=targetgreen!80!black] at (3.0, 0.6) {$(1, y)$};
    \node[left, font=\scriptsize, text=targetgreen!80!black] at (0.2, 0.6) {$(0, y)$};
    
    \node[below, font=\footnotesize, align=center, text=gray!80!black] at (1.6, -1.3) 
      {Roll into cylinder, then\\join circular open ends.};
  \end{scope}

  % PANEL 3: Möbius Strip (Half-twist)
  \begin{scope}[shift={(9.6,0)}]
    \node[font=\small\bfseries, text=trapred!80!black] at (1.6, 2.6) {Mobius Band: Half-Twist};
    
    % Square
    \fill[quotientorange!5] (0.2, -0.8) rectangle (3.0, 2.0);
    \draw[gray!30, thick] (0.2, -0.8) rectangle (3.0, 2.0);
    
    % Left edge pointing UP
    \draw[edgestyleb, -{Latex[length=3mm]}] (0.2, -0.8) -- (0.2, 0.8);
    \draw[edgestyleb] (0.2, 0.8) -- (0.2, 2.0);
    \node[left, font=\scriptsize, text=targetgreen!80!black] at (0.2, 0.6) {$(0, y)$};
    
    % Right edge pointing DOWN
    \draw[edgestylerev, -{Latex[length=3mm]}] (3.0, 2.0) -- (3.0, 0.4);
    \draw[edgestylerev] (3.0, 0.4) -- (3.0, -0.8);
    \node[right, font=\scriptsize, text=trapred!80!black] at (3.0, 0.6) {$(1, 1-y)$};
    
    \node[below, font=\footnotesize, align=center, text=gray!80!black] at (1.6, -1.3) 
      {Reversed vertical arrows\\force a $180^\circ$ flip: one-sided.};
  \end{scope}
\end{tikzpicture}
\caption{Geometric gluing via equivalence relations. \textbf{Left}: Gluing the endpoints of the unit interval $[0, 1]$ via $0 \sim 1$ creates the circle $S^1$. \textbf{Middle}: Identifying opposite edges of a square with matching orientations ($(0, y) \sim (1, y)$ and $(x, 0) \sim (x, 1)$) first rolls the square into a cylinder and then joins the two circular boundaries to form the torus $T^2$ (a doughnut surface). \textbf{Right}: Gluing the vertical edges with a half-twist ($(0, y) \sim (1, 1-y)$, shown by opposite arrowheads) creates the Möbius strip, a non-orientable surface with only one side and a single boundary loop.}
\label{33:fig-geometric-gluing}
\end{figure}

As illustrated in Figure~\ref{33:fig-geometric-gluing}, three classical spaces emerge from simple identifications:
\begin{itemize}
  \item \textbf{The Circle $S^1$}: Take the unit interval $[0, 1] \subset \R$. Identify only the two boundary endpoints: $0 \sim 1$, leaving all interior points $t \sim t$ alone. Identifying $0$ and $1$ joins the endpoints to form a closed loop: the circle $S^1$.
  \item \textbf{The Torus $T^2$}: Take the flat unit square $[0, 1] \times [0, 1]$. Glue the left edge to the right edge with matching orientation ($(0, y) \sim (1, y)$), rolling the square into a hollow tube or cylinder. Then bend the tube and glue the bottom circular boundary to the top circular boundary ($(x, 0) \sim (x, 1)$). The resulting quotient is the 2-dimensional torus $T^2$, the surface of a doughnut.
  \item \textbf{The Möbius Strip}: Take $[0, 1] \times [0, 1]$ and glue the vertical ends with a half-twist: $(0, y) \sim (1, 1-y)$. The reversed arrowheads force the surface to flip upside down before reattaching. The resulting quotient is \emph{non-orientable}: traversing the center line once returns to the starting point with reversed orientation, demonstrating that the surface has only a single continuous side and a single boundary loop.
\end{itemize}

\subsection{What Quotients Unlock---and Their Limits}

The invention of quotients enabled several major developments across mathematics:
\begin{enumerate}
  \item \textbf{Construction of the Continuum}: The real numbers $\R$ can be constructed with complete logical rigor from the rational numbers $\Q$. A Cauchy sequence is a sequence of fractions whose terms become arbitrarily close; a null sequence is one that converges to zero. By declaring two Cauchy sequences equivalent whenever their difference is a null sequence, the quotient set of Cauchy sequences modulo null sequences completes $\Q$, constructing the continuous real line $\R$.
  \item \textbf{Topological Surgery}: Manifolds---geometric spaces that locally resemble flat Euclidean space $\R^n$, just as a sphere locally resembles a flat plane---are systematically constructed by gluing polyhedral building blocks along their faces. For example, the real projective plane $\mathbb{RP}^2$ is built by identifying antipodal points ($p \sim -p$) on the boundary of a hemisphere.
  \item \textbf{Decomposition of Symmetries}: In group theory, quotient groups $G/N$ allow algebraists to factor complicated transformation groups into simple atomic pieces (groups with no non-trivial quotients). This structural factoring forms the backbone of the Jordan--H\"older program, which completely classified all finite simple groups.
\end{enumerate}

\paragraph{Scope and Limitations.}
Yet quotient constructions carry serious hazards:
\begin{itemize}
  \item \textbf{Pathological Quotient Topologies}: Gluing points can destroy geometric stability. A quotient of a smooth, well-behaved manifold can fail to be Hausdorff (distinct points cannot be separated by disjoint open sets) or fail to be locally Euclidean, creating sharp cone points, self-intersecting cusps, or orbifolds (spaces that look locally like flat Euclidean space divided by a finite group of rotations).
  \item \textbf{Representative Overhead and Non-Measurable Sets}: In an infinite quotient space, there is often no continuous or algorithmic rule to select a single representative from each class. The Axiom of Choice is required just to pick a set of representatives. In the quotient group $\R/\Q$, using the Axiom of Choice to pick one representative from each coset in $[0, 1]$ constructs the Vitali set---a subset of the real line so fragmented that it cannot be assigned any consistent length or probability.
\end{itemize}

\paragraph{Where We Go Next in the Book.}
Having examined how equivalence relations partition sets and how quotients collapse structure, we move from static spaces to inductive hierarchies. In Chapter~34 (\emph{Second Pass: Induction Beyond Formulas}), we leave behind routine algebraic summations to uncover the Well-Ordering Principle and structural induction, proving theorems across recursive trees, compilers, and ordered structures.

\subsection{Exercises}

\begin{enumerate}
  \item \textbf{Conceptual and Mental Warm-ups}:
  \begin{enumerate}
    \item If $\sim$ is an equivalence relation on a non-empty set $A$, can an equivalence class $[a]$ ever be empty? Explain using reflexivity.
    \item If $a \sim b$, what is $[a] \cap [b]$? What is $[a] \cup [b]$?
    \item Give a concrete relation on $\{1, 2, 3\}$ that is reflexive and symmetric, but fails transitivity.
  \end{enumerate}

  \item \textbf{Concentric Circles in $\R^2$}:
  Let $\sim$ be the relation on $\R^2$ defined by $(x_1, y_1) \sim (x_2, y_2)$ if and only if $x_1^2 + y_1^2 = x_2^2 + y_2^2$.
  \begin{enumerate}
    \item Verify that $\sim$ is an equivalence relation on $\R^2$.
    \item Describe the equivalence class $[(0, 0)]$. For any $(x_0, y_0) \neq (0, 0)$, what geometric shape is $[(x_0, y_0)]$?
    \item Consider the radius map $r: \R^2 \to [0, \infty)$ defined by $r(x, y) = \sqrt{x^2 + y^2}$. Show that $\sim$ is the fiber relation of $r$, and deduce that $\R^2/\sim \;\cong\; [0, \infty)$ as sets.
  \end{enumerate}

  \item \textbf{Modular Congruence on $\Z$}:
  Define a relation on $\Z$ by $a \sim b \iff 4 \mid (a - b)$.
  \begin{enumerate}
    \item Prove that $\sim$ is an equivalence relation.
    \item Write out the four distinct equivalence classes $[0], [1], [2], [3]$ using set notation.
    \item Prove that class addition $[a] + [b] = [a+b]$ and class multiplication $[a] \cdot [b] = [ab]$ are well-defined: if $a \sim a'$ and $b \sim b'$, prove that $(a+b) \sim (a'+b')$ and $(ab) \sim (a'b')$.
  \end{enumerate}

  \item \textbf{The Exponential Map and the Circle}:
  Let $f: \R \to \C^\times$ be the map $f(t) = e^{2\pi i t} = \cos(2\pi t) + i \sin(2\pi t)$.
  \begin{enumerate}
    \item Identify the image $\operatorname{im}(f) \subset \C^\times$.
    \item Show that $f(s) = f(t)$ if and only if $s - t \in \Z$.
    \item Apply the First Isomorphism Theorem for Sets to conclude that $\R/\Z \cong S^1$.
  \end{enumerate}

  \item \textbf{Constructing the Rational Numbers $\Q$ (Challenge)}:
  Let $S = \Z \times (\Z \setminus \{0\})$, the set of pairs of integers $(p, q)$ with nonzero denominator $q$. Define a relation on $S$ by:
  \[
  (p, q) \sim (r, s) \iff ps = qr.
  \]
  \begin{enumerate}
    \item Prove that $\sim$ is an equivalence relation on $S$. (Be careful: where exactly does your proof of transitivity use the fact that denominators are nonzero?)
    \item Define addition on equivalence classes by $[(p, q)] + [(r, s)] = [(ps + qr, qs)]$. Prove that this operation is well-defined.
    \item Why does the formula $N([(p, q)]) = p$ fail to define a function on $S/\sim$?
  \end{enumerate}

  \item \textbf{The Real Projective Line $\mathbb{RP}^1$ (Geometry Challenge)}:
  On the punctured plane $V = \R^2 \setminus \{(0, 0)\}$, define two nonzero vectors to be equivalent, $v_1 \sim v_2$, if there exists a nonzero scalar $\lambda \in \R^\times$ such that $v_1 = \lambda v_2$.
  \begin{enumerate}
    \item Prove that $\sim$ is an equivalence relation.
    \item Show that each equivalence class $[(x, y)]$ is an entire straight line passing through the origin (minus the origin itself).
    \item Show that each equivalence class intersects the unit circle $S^1 = \{(x, y) \in \R^2 : x^2 + y^2 = 1\}$ in exactly two antipodal points: $(x, y)$ and $(-x, -y)$. Conclude that $\mathbb{RP}^1 \cong S^1/(p \sim -p)$.
  \end{enumerate}
\end{enumerate}

\newpage
